\scanpage{178}

arriver en tous cas à définir une équivalence \uuline{canonique} sur les groupoïdes fondamentaux des $\unsure{\check N|\breve N}^{*}_{d_*;(d'_0,d''_0)}|D_{d_*}^{*}$ d'une part, des $\unsure{V}_{d_*,d''_0}^{d'_0}$ d'autre part
\[
(\unsure{295})\qquad
\begin{array}{|l}
\hline
\displaystyle
\Pi_1\left(\underbrace{\unsure{\check N}^{*}_{d_*;(d'_0,d''_0)}|D_{d_*}^{*}}_{\overset{\mathrm{def}}{=}\,\unsure{\check N}^{**}_{d_*;(d'_0,d''_0)}}\right)
\xrightarrow{\approx}\Pi_1\unsure{V}_{d_*,d''_0}^{d'_0}\\
\hline
\end{array}
\]

Admettant ce point, la connaissance de $\Pi_1\unsure{V}_{d_*,d''_0}^{d'_0}$ \deleted{\unsure{revient donc}} \inserted{est équivalente} à celle des $\pi_1$ ordinaires de $\unsure{\check N}^{**}$ \deleted{\unsure{qui}} au dessus des différentes composantes connexes des $D_{d_*}^{*}$ --- et ce sont, par la suite exacte d'homotopie, des extensions des $\pi_1$ des comp. conn. des $D_{d_*}^{*}$, par des quotients de $\mathbb Z^{\unsure{r|v}}$ (par l'image des $\pi_2$ desdites comp. connexes\ldots)

\margin{177}
